\documentclass[utf8]{ctexart}
\usepackage{graphicx}
\usepackage{amsmath, amssymb, amsthm}
\usepackage{geometry}
\usepackage{enumitem}
\usepackage{xcolor}
\usepackage{mdframed}
\usepackage{tikz}

\usetikzlibrary{arrows.meta,positioning,calc}
\usepackage{booktabs}
\usepackage{array}

\geometry{margin=2cm}

\definecolor{critical}{RGB}{200,40,40}
\definecolor{nodeblue}{RGB}{52,120,200}
\definecolor{nodegray}{RGB}{215,228,242}
\definecolor{critnodebg}{RGB}{255,218,218}

\usetikzlibrary{graphs, positioning, arrows.meta}

\geometry{a4paper, margin=2.5cm}

% 定理环境
\newtheoremstyle{mystyle}
  {8pt}{8pt}{\kaishu}{0pt}{\bfseries}{}{1em}{}
\theoremstyle{mystyle}
\newtheorem{theorem}{定理}
\newtheorem{lemma}{引理}

% 好看的题目框
\newmdenv[
  linecolor=black!40,
  linewidth=1pt,
  backgroundcolor=gray!8,
  roundcorner=5pt,
  innertopmargin=8pt,
  innerbottommargin=8pt,
]{problembox}

\title{\textbf{离散数学2\quad 作业7}}
\date{\today}
\author{李子嘉 2024012325}
\begin{document}
\maketitle

\section*{题目1}

\begin{problembox}
    设 $ G = (V, E) $ 是一个有向连通图，$ e $ 是 $ G $ 中的一条边，证明：
    \begin{enumerate}
        \item 若 $ e $ 不在 $ G $ 的任何一棵支撑树中，则 $ e $ 为自环。
        \item 若 $ e $ 在 $ G $ 的每一棵支撑树中，则 $ e $ 为割边。
    \end{enumerate}
\end{problembox}
\begin{proof}
我们先证明第一问：
假设 $e$ 不是自环，那么 $e$ 连接了两个不同的顶点 $u$ 和 $v$。
我们讨论两种情况：
\begin{enumerate}
    \item 如果 $e$ 是割边，那么删除 $e$ 后图就不连通了，形成了两个连通分量，$u$ 和 $v$ 分别位于不同的分量中。而支撑树必须连接所有顶点，必须连接 $u$ 和 $v$，因此 $e$ 必须在支撑树中，矛盾。
    \item 如果 $e$ 不是割边，那么删除 $e$ 后图仍然是连通的。我们可以在删除后的图 $G'$ 中找到一棵支撑树 $T$ ，这棵支撑树在原图 $G$ 中也是支撑树，并且不包含边 $e$。
    由于 $u$ 和 $v$ 仍在图 $G'$ 中，所以$T$ 中也连接了 $u$ 和 $v$，但不经过边 $e$，我们记录这一路径上的边的集合，设为 $P$。
    现在，我们把边 $e$ 加入 $T$ 中，得到 $T'$，这个图不是树，因为有一个环 $C$，这个环包含边 $e$ 和路径 $P$ 上的边。
    但，我们可以从 $P$ 中删除任意一条边，得到一个新的树 $T''$，这个树仍然连接了所有顶点，并且包含边 $e$，这与题设矛盾。
\end{enumerate}

因此，$e$ 必须是自环。

接下来，我们用反证法证明第二问：

假设 $ e $ 不是割边，即删除 $ e $ 后图 $ G' $ 仍然是连通的，并且顶点集与 $ G $ 相同。
由于 $ G' $ 是连通的，我们可以找到一棵支撑树 $ T $ ，这棵支撑树在原图 $ G $ 中也是支撑树（因为 $ T $ 同样连接了所有顶点）。
但由于我们是在 $G'$ 里面找的 $T$ ，它肯定不包含边 $ e $，所以 $e$ 并不在 $ G $ 的每一棵支撑树中，这与题设矛盾，因此 $ e $ 必须是割边。
\end{proof}

\newpage
\section*{题目2}
\begin{problembox}
    求图3.26所示有向图中必含(v1, v2)、(v2, v7)、(v4, v6)的支撑树的数目。
\end{problembox}
\textbf{解：} 题目钦定了边 $(v1, v2)$、$(v2, v7)$、$(v4, v6)$ 必须在支撑树中，
我们就得合并1、2、7为一个超级顶点 $v_{127}$，4、6为一个超级顶点 $v_{46}$。
我们更新关联矩阵，得到一个新的图 $G'$，在这个图中求支撑树的数目。我们给边标序号：
\medskip

\begin{tabular}{|c|c|c|c|c|c|c|c|c|c|}
\hline
1 & 2 & 3 & 4 & 5 & 6 & 7 & 8 & 9 \\
\hline
$(v_{127}, v_{46})$ & $(v_{127}, v_{46})$ & $(v_{127}, v_{3})$ & $(v_{127}, v_{3})$ & $(v_{127}, v_{5})$ & $(v_{5}, v_{127})$ & $(v_{5}, v_{46})$ & $(v_{5}, v_{46})$ & $(v_{3}, v_{5})$ \\
\hline
\end{tabular}
\medskip

顶点序按照 $v_{127}, v_{3}, v_{46}, v_{5}$ 的顺序排列，得到关联矩阵 $A$：
\[
A = \begin{bmatrix}
1 & 1 & 1 & 1 & 1 & -1 & 0 & 0 & 0 \\
0 & 0 & -1 & -1 & 0 & 0 & 0 & 0 & 1 \\
-1 & -1 & 0 & 0 & 0 & 0 & -1 & -1 & 0 \\
0 & 0 & 0 & 0 & -1 & 1 & 1 & 1 & -1 \\
\end{bmatrix}
\]
我们删除最复杂的第1行，得到矩阵 $A'$：
\[
A' = \begin{bmatrix}
0 & 0 & -1 & -1 & 0 & 0 & 0 & 0 & 1 \\
-1 & -1 & 0 & 0 & 0 & 0 & -1 & -1 & 0 \\
0 & 0 & 0 & 0 & -1 & 1 & 1 & 1 & -1 \\
\end{bmatrix}
\]
我们计算 $A'A'^T$ 的行列式，得到 $A'A'^T$ 为：
\[
A'A'^T = \begin{bmatrix}3 & 0 & -1 \\ 0 & 4 & -2 \\ -1 & -2 & 5\end{bmatrix}
\]

它的行列式为 $44$，因此图 $G'$ 的支撑树数目为 $44$。

\newpage
\section*{题目3}
\begin{problembox}
在图3.28中，求：
\begin{enumerate}
    \item 以v1为根的根树的数目。
    \item 以v1为根，不含(v1, v5)的根树的数目。
    \item 以v1为根，包含(v2, v3)的根树的数目。
\end{enumerate}
\end{problembox}
\textbf{解：} 我们先给出图3.28的关联矩阵 $A$，按照顶点 $v1, v2, v3, v4, v5$ 的顺序排列。边的序号为：
\begin{center}

\begin{tabular}{|c|c|}
\hline
编号 & (起点, 终点) \\
\hline
1 & $(v1, v2)$ \\
2 & $(v1, v4)$ \\
3 & $(v1, v5)$ \\
4 & $(v5, v2)$ \\
5 & $(v2, v5)$ \\
6 & $(v5, v4)$ \\
7 & $(v2, v3)$ \\
8 & $(v5, v3)$ \\
9 & $(v3, v4)$ \\
10 & $(v4, v3)$ \\
\hline
\end{tabular}
\end{center}

\[
A = \begin{bmatrix}
1 & 1 & 1 & 0 & 0 & 0 & 0 & 0 & 0 & 0 \\
-1 & 0 & 0 & -1 & 1 & 0 & 1 & 0 & 0 & 0 \\
0 & 0 & 0 & 0 & 0 & 0 & -1 & -1 & 1 & -1 \\
0 & -1 & 0 & 0 & 0 & -1 & 0 & 0 & -1 & 1 \\
0 & 0 & -1 & 1 & -1 & 1 & 0 & 1 & 0 & 0 \\
\end{bmatrix}
\]
\subsection*{第一问}
以1为根，我们需要删除第1行，随后把所有1元素改为0，得到矩阵 $A'$：
\[
A' = \begin{bmatrix}
-1 & 0 & 0 & -1 & 0 & 0 & 0 & 0 & 0 & 0 \\
0 & 0 & 0 & 0 & 0 & 0 & -1 & -1 & 0 & -1 \\
0 & -1 & 0 & 0 & 0 & -1 & 0 & 0 & -1 & 0 \\
0 & 0 & -1 & 0 & -1 & 0 & 0 & 0 & 0 & 0 \\
\end{bmatrix}
\]
我们计算 $A'A'^T$ 的行列式，得到 $A'A'^T$ 为：
\[
A'A'^T = \begin{bmatrix}
2 & 0 & 0 & -1 \\
-1 & 3 & -1 & -1 \\
0 & -1 & 3 & -1 \\
-1 & 0 & 0 & 2
\end{bmatrix}
\]
它的行列式为 $24$，因此以 $v1$ 为根的根树数目为 $24$。
\subsection*{第二问}
以 $v1$ 为根，不含 $(v1, v5)$ 的根树的数目，我们需要在 $A'$ 的基础上把第3列删掉，得到矩阵 $A''$：
\[
A'' = \begin{bmatrix}
-1 & 0 & -1 & 0 & 0 & 0 & 0 & 0 & 0 \\
0 & 0 & 0 & 0 & 0 & -1 & -1 & 0 & -1 \\
0 & -1 & 0 & 0 & -1 & 0 & 0 & -1 & 0 \\
0 & 0 & 0 & -1 & 0 & 0 & 0 & 0 & 0 \\
\end{bmatrix}
\]
类似地，我们计算 $A''A''^T$ 的行列式，得到 $A''A''^T$ 为：
\[
A''A''^T = \begin{bmatrix}2 & 0 & 0 & -1 \\ -1 & 3 & -1 & -1 \\ 0 & -1 & 3 & -1 \\ -1 & 0 & 0 & 1\end{bmatrix}
\]
它的行列式为 $8$，因此以 $v1$ 为根，不含 $(v1, v5)$ 的根树数目为 $8$。

\subsection*{第三问}
以 $v1$ 为根，包含 $(v2, v3)$ 的根树的数目，我们需要在 $A'$ 的基础上把第8列和第10列都删掉，得到矩阵 $A'''$：
\[
A''' = \begin{bmatrix}
-1 & 0 & 0 & -1 & 0 & 0 & 0 & 0 \\
0 & 0 & 0 & 0 & 0 & 0 & -1 & 0 \\
0 & -1 & 0 & 0 & 0 & -1 & 0 & -1 \\
0 & 0 & -1 & 0 & -1 & 0 & 0 & 0 \\
\end{bmatrix}
\]
类似地，我们计算 $A'''A'''^T$ 的行列式，得到 $A'''A'''^T$ 为：
\[
A'''A'''^T = \begin{bmatrix}
2 & 0 & 0 & -1 \\
-1 & 1 & 0 & 0 \\
0 & -1 & 3 & -1 \\
-1 & 0 & 0 & 2
\end{bmatrix}
\]
它的行列式为 $9$，因此以 $v1$ 为根，包含 $(v2, v3)$ 的根树数目为 $9$。

\section*{附录}
本次作业由于涉及到大量的矩阵运算，手算或计算器已经无法胜任，因此我编写了一个交互式的矩阵计算器来辅助完成作业。这个工具支持输入、显示、加法、乘法、转置、行列式、删除行列等功能，可以在命令行中运行，输入矩阵时支持纯文本和LaTeX格式，输出结果也可以显示为纯文本或LaTeX格式，方便复制到作业中。
具体代码由 GPT-5.3-Codex 实现，后续我对它进行了一些修改和完善。代码文件名为 \texttt{MatrixCalculator.py}，在作业的同一目录下，可以直接运行使用。
值得注意的是，此工具仅代替了计算器的功能，作业中所有的解题思路和计算都是我独立完成的，工具只是为了提高效率和准确性。
另外，第二题的两种算法实现（合并节点法和容斥法）以及第三题的删除法和容斥法的结果也都在 \texttt{spanning\_tree.py} 文件中进行了展示和比较，结果都是一致的，验证了算法的正确性。

\end{document}